\documentclass[UTF8,12pt,a4paper,fontset=fandol]{ctexbook}


% 支持从项目根目录或当前笔记目录编译。
\makeatletter
\def\input@path{{../}}
\makeatother
\usepackage{researchnotes}




\begin{document}

\chapter{函数}
数学分析主要由微积分和级数理论组成，它所研究的主要对象是实函数，即以实数为自变量并且在实数中取值的函数。因此，在本章中我们首先简要地介绍一下实数系的连续性，然后介绍函数的概念和有关的基本知识。


\section{实数}
在近几个世纪中，科学技术之所以取得了辉煌的成就，在很大程度上是因为数学研究取得了重大进展，其中微积分的创立是现代数学的里程碑。微积分在物理、天文、技术、化学、生物等的研究中显示了强大的威力，解决了许多过去认为高不可攀的困难问题，促进了科学技术的发展。然而，由于在创建初期微积分是以几何直观和物理直觉为依据而进行演绎推理的，因此就形成了方法上有效但逻辑上不能自圆其说的矛盾局面。为了解决微积分在理论上面临的问题，许多著名数学家都投身于微积分理论基础的研究。人们后来发现，微积分的主要理论基础是严格的极限理论。到了 19 世纪初，柯西（Cauchy）以极限理论为微积分奠定了理论基础。但是柯西构筑的理论大厦起初并不完善，这是因为柯西并没有对实数给出严格的定义。而后来人们又发现，极限理论的某些基本原理依赖于实数系的连续性。为此，本节简要地介绍一下这方面的内容。


\subsection{数集}
在介绍实数理论前，我们简要地介绍一下集合的概念。在数学中，所谓\impo{集合}是将具有某种特性的对象的全体放在一起作为一个整体，通常用大写字母 $A,\,B,\,C,\,X,\,Y$ 等记之。这些对象称为集合中的\impo{元素}，通常用小写字母 $a,\,b,\,c,\,x,\,y$ 等记之。若 $a$ 是 $A$ 中的元素，则记为 $a\in A$，并读做 $a$ 属于 $A$。如果 $a$ 不是 $A$ 中的元素，则记为 $a\notin A$，并读做 $a$ 不属于 $A$。集合的概念是最基本的概念，它不能用别的概念来加以定义。


集合通常有列举法和描述法两种表示法。列举法是将集合中的元素全部列出，如：$A=\{1,\,2,\,3,\,\cdots,\,10\}$；而描述法是将集合的特性精确给出，如：$B=\{x:x\,\text{是}\,x^2-1=0\,\text{的根}\}$；$E=\{x:x\,\text{是正有理数}\}$。一般来说，在数学分析中我们主要研究由数构成的集合，并且很少研究由有限个元素构成的集合；大部分集合是由无穷多个元素组成 (这种集合简称为无穷集合)，因此这些集合必须用描述法表示。


若集合 $A$ 中的每一个元素 $x$ 都属于集合 $B$，则称 $B$ 包含 $A$，记为 $A \subseteq B$，此时也称 $A$ 是 $B$ 的\impo{子集}。如果 $A \subseteq B$ 和 $B \subseteq A$ 同时成立，则认为 $A,\,B$ 是同一个集合，此时也记为 $A=B$。如上面的集合 $B=\{x:x\, \text{是}\, x^2-1=0\,\text{的根}\}$ 与 $C=\{1,\, -1\}$ 是相等的。


若 $A\subseteq B$ 且 $A\ne B$，则称 $A$ 是 $B$ 的\impo{真子集}，记为 $A\subset B$。特别地，我们引入\impo{空集} $\varnothing$，即 $\varnothing$ 中不含有任何元素，因此 $\varnothing$ 是任何集合的子集。


集合有如下的运算：给定集合 $A,\, B$，有
\begin{itemize}
  \item $A\cup B=\{x:x\in A \,\text{或}\, x\in B\}$ 称为 $A$ 与 $B$ 的\impo{并}；
  \item $A\cap B=\{x:x\in A\,\text{且}\,x\in B\}$ 称为 $A$ 与 $B$ 的\impo{交}；
  \item $A\setminus B=\{x:x\in A\,\text{且}\,x\notin B\}$ 称为 $A$ 与 $B$ 的\impo{差}。
\end{itemize}


为了数学论述简便，习惯上常将逻辑推理中一些常用的词用符号表示。现将本书常用的两个数学记号 $\forall$ 和 $\exists$ 说明如下：$\forall$ 代表“任意一个”；$\exists$ 代表“存在”。利用上述记号可以使得我们以后在行文上非常方便。例如，我们可以将 $A\subset B$ 的定义叙述为：$\forall x\in A$，有 $x\in B$。


\subsection{实数系的连续性}
人类所认识的第一个数系就是自然数系，它的定义为 $\{0,\,1,\,2,\,3,\,\cdots\}$。虽然自然数对于计数来说是够用了，但是自然数系并不是一个完善的数系。首先作为量的描述手段，它只能表示一个单位量的整数倍，而无法去表示此单位量的部分。此外，作为量的运算手段，它只能自由地进行加、乘运算，而不能自由地进行加、乘运算的逆运算。自然数的这种离散性和运算的不完备性，促使人们去对它进行扩充。人们首先引进了负数，得到了整数系。对整数系人们可以自由进行加、减运算。人们为了得到可以自由地进行加、减、乘、除四则运算的数系，对整数系中的任意两个数进行加、减、乘、除（除数不为零）得到的数的全体记为一个新的数系，这就是随后得到的\impo{有理数系}。


对于一个数集 $K$，若 $K$ 中至少有一个非零元素，且 $K$ 中任何两个元素的加、减、乘、除（除数不为零）运算后得到的数仍然属于 $K$，即 $K$ 关于四则运算封闭，则称 $K$ 为一个数域。容易看出，有理数集就是一个数域。从代数上来讲，有理数系已经是一个完美的数系，它可以在任何精度要求下对一个量进行表示和实施有效的运算，并且任何两个有理数之间必有有理数存在（或说有理数有稠密性）。然而，早在 2500 年前，人们就发现有理数系也有缺陷。例如，若用 $c$ 来表示一个边长为 $1$ 的正方形的对角线的长度，则 $c$ 就无法用有理数表示。所以，有理数虽然在数轴上密密麻麻，但并没有布满整个数轴，留有许多“空隙”。这说明还有新的数存在，但它没有被严格地定义。


由于有理数系已经对四则运算封闭，因此我们必须用不同于以前的数系的扩充才有可能得到新的数。我们前面曾经指出，微积分的理论基础是极限论，但在这种具有“空隙”的有理数系上实施极限运算，就会极为不便。因此，正如四则运算需要一个封闭的数域一样，极限运算也需要一个关于极限封闭的数域。这就需要将有理数进行扩充，使其能够填补有理数在数轴上留下的所有这些“空隙”。于是，这又归结到如何定义无理数的问题上了。这一历史任务，终于在 19 世纪后半叶，由戴德金（Dedekind）和康托尔（Cantor）等人完成。以下我们简要地介绍一下戴德金关于实数的构造方法。


\begin{definition}
设 $S$ 是一个有大小顺序的非空数集，$A$ 和 $B$ 是它的两个子集，如果它们满足以下条件：
\begin{enumerate}[label=(\arabic*)]
  \item $A\ne\varnothing$，$B\ne\varnothing$；
  \item $A\cup B=S$；
  \item $\forall a\in A$，$\forall b\in B$，都有 $a<b$；
  \item $A$ 中无最大数，
\end{enumerate}
则我们将 $A,\,B$ 称为 $S$ 的一个\impo{分划}，记为 $(A\mid B)$。
\end{definition}


现在我们来考虑有理数系 $\mathbb{Q}$ 的分划。对 $\mathbb{Q}$ 的任意分划 $(A\mid B)$，必有以下两种情形之一发生：
\begin{enumerate}[label=(\arabic*)]
  \item $B$ 中存在最小数，此时称 $(A\mid B)$ 是一个\impo{有理分划}；
  \item $B$ 中不存在最小数，此时称 $(A\mid B)$ 是一个\impo{无理分划}。
\end{enumerate}


\begin{example}
记
\[
  A=\{x\in\mathbb{Q}:x\leqslant 0\,\text{或}\, x>0\,\text{且}\,x^2<2\},\qquad B=\mathbb{Q}\setminus A,
\]
则 $(A\mid B)$ 是一个无理分划。
\end{example}


有理数系 $\mathbb{Q}$ 的所有分划构成了一个集合，我们称这个集合为\impo{实数系}，并记为 $\mathbb{R}$。显然，有理数集与 $\mathbb{R}$ 中的有理分划是一一对应的。因此 $\mathbb{R}$ 可以被认为是由有理数集加上无理分划所构成。我们称 $\mathbb{R}$ 中的这些无理分划为\impo{无理数}。换句话说，$\mathbb{R}$ 是由全体有理数与无理数所构成的集合。例如，例 1.1.1 中的无理分划对应的就是大家熟知的无理数 $\sqrt{2}$。我们可以在 $\mathbb{R}$ 上很自然地定义大小顺序关系及四则运算，经过冗长但不是很困难的论证，可以证明它是一个以有理数域为其子域的有序域。


前面我们已经指出有理数系有稠密性，但没有连续性，即有理数之间有许多“空隙”。下面的戴德金分割定理告诉我们，在有理数集合中加入无理数之后，就没有“空隙”了，也就是说实数系具有了连续性。对于一个数集 $A$，若对任意两个数 $a,\,b\in A$，$a,\,b$ 之间的所有的数都在 $A$ 中，则称 $A$ 是\impo{连通的}。换句话说，实数集是一个连通的集合。


\begin{theorem}
（戴德金分割定理）对 $\mathbb{R}$ 的任一分划 $(A\mid B)$，$B$ 中必有最小数。
\end{theorem}


定理 1.1.1 告诉我们，对 $\mathbb{R}$ 再进行分划就不可能产生新数了。换句话说，$\mathbb{R}$ 已将整个实轴填满了。因此 $\mathbb{R}$ 是一个有序的连通域。对于每一个实数 $x$，若 $x\geqslant 0$，则它在原点 $O$ 的右边且它到 $O$ 的距离为 $x$；若 $x<0$，则它在原点 $O$ 的左边且它到 $O$ 的距离为 $-x$。根据上面的对应法则，显然，对于 $\mathbb{R}$ 中的每个数，在数轴上有且只有一个点与之对应。反之，任给数轴上的一个点 $A$，若 $A$ 就是原点 $O$，则它与 $0\in\mathbb{R}$ 对应；若 $A$ 在原点 $O$ 的右边，则它与 $A$ 到 $O$ 的距离 $x\in\mathbb{R}$ 对应；若 $A$ 在原点 $O$ 的左边，则它与 $-x$ 对应，其中 $x$ 为 $A$ 到 $O$ 的距离。这样，我们就建立了 $\mathbb{R}$ 中的数与数轴上的点之间的一一对应。因此，今后我们将不再区分数轴上的点与它所对应的实数。


\subsection{有界集与确界}
设集合 $E\subseteq\mathbb{R}$，并且 $E\ne\varnothing$。如果存在 $M\in\mathbb{R}$，使得对 $\forall x\in E$，有 $x\leqslant M$，则称 $E$ 是\impo{有上界的}，并且说 $M$ 是 $E$ 的一个\impo{上界}；如果存在 $m\in\mathbb{R}$，使得对 $\forall x\in E$，有 $x\geqslant m$，则称 $E$ 是\impo{有下界的}，并且说 $m$ 是 $E$ 的一个\impo{下界}；如果 $E$ 既有上界又有下界，则称 $E$ 是\impo{有界的}。显然，$E$ 是有界的充分必要条件是：存在 $M>0$，使得对任意的 $x\in E$，有 $|x|\leqslant M$。


\begin{example}
集合
\[
  E=\left\{0,\,\frac32,\,\frac23,\,\cdots,\,1+\frac{(-1)^n}{n},\,\cdots\right\}
\]
是有界集，它的一个上界是 $\dfrac32$，一个下界是 $0$。
\end{example}


\begin{example}
集合
\[
  E=\left\{1,\,\frac12,\,\frac13,\,\cdots,\,\frac1n,\,\cdots\right\}
\]
是有界集，其中 $1$ 是它的一个上界，$0$ 是它的一个下界。
\end{example}


若 $E$ 是一个有界数集时，它的上（下）界必有无穷多个。是否有一个最小（大）的上（下）界呢？为了对此进行精确描述，我们引入：


\begin{definition}
设 $E\subset\mathbb{R}$ 为一个非空数集。若有 $M\in\mathbb{R}$ 满足
\begin{enumerate}[label=(\arabic*)]
  \item $M$ 是 $E$ 的一个上界，即 $\forall x\in E$，有 $x\le M$；
  \item 对 $\forall\varepsilon>0$，存在 $x'\in E$，使得 $x'>M-\varepsilon$，
\end{enumerate}
则称 $M$ 为 $E$ 的\impo{上确界}，记为 $M=\sup E=\sup_{x\in E}\{x\}$。


若有 $m\in\mathbb{R}$ 满足
\begin{enumerate}[label=(\arabic*)]
  \item $m$ 是 $E$ 的一个下界，即 $\forall x\in E$，有 $x\ge m$；
  \item 对 $\forall\varepsilon>0$，存在 $x'\in E$，使得 $x'<m+\varepsilon$，
\end{enumerate}
则称 $m$ 为 $E$ 的\impo{下确界}，记为 $m=\inf E=\inf_{x\in E}\{x\}$。
\end{definition}

从定义容易看出，$E$ 的上确界就是它的最小上界，而下确界就是它的最大下界。如果 $\sup E\in E$，则 $\sup E$ 可记为 $\max E$，这时上确界即为 $E$ 中的最大数；如果 $\inf E\in E$，则 $\inf E$ 可记为 $\min E$，这时下确界即为 $E$ 中的最小数。


此外，引入记号 $+\infty$ 和 $-\infty$，分别表示正无穷大和负无穷大，简称正无穷和负无穷。为了以后行文方便，我们约定：当数集 $E$ 无上界时，记做 $\sup E=+\infty$；当数集 $E$ 无下界时，记做 $\inf E=-\infty$。这里需特别指出的是，$+\infty$ 和 $-\infty$ 只是记号而不是实数。当数集 $E$ 满足 $\sup E=+\infty$（$\inf E=-\infty$） （即 $E$ 是无上界（下界））时，它的等价说法是：$\forall M>0$，$\exists x\in E$，并且有 $x>M$（$x<-M$）。


\begin{example}
设
\[
  E=\left\{\frac pq:p,q\,\text{为正整数，而且满足}\,p\leqslant q\right\},
\]
则有 $\inf E=0\notin E,\ \sup E=1=\max E$。这表明，$E$ 中没有最小数，而有最大数 $1$。
\end{example}


下面的定理与戴德金分割定理等价，即如下定理是实数系连续性的另一种表述，但它使用起来特别方便。


\begin{theorem}
  （确界存在定理）非空有上界的实数集必有上确界；非空有下界的实数集必有下确界。
\end{theorem}


\noindent\impo{证明}\quad 我们只证明上确界的情形，下确界的情形完全类似。设 $E$ 是一个非空有上界的实数集。若 $E$ 中存在最大数 $M$ 时，则
\[
  \sup E=\max E=M.
\]
现假设 $E$ 中没有最大数。我们对 $\mathbb{R}$ 作分划：$B$ 是由 $E$ 的所有上界组成的集合，而 $A=\mathbb{R}\setminus B$。由 $E$ 的有界性，推出 $B\ne\varnothing$；而由 $E\ne\varnothing$，推出 $A\ne\varnothing$。显然，对任意 $a\in A$ 和 $b\in B$，有 $a<b$，而且 $A$ 中无最大数。因此，$(A|B)$ 是 $\mathbb{R}$ 的一个分划，从而由戴德金分割定理知，$B$ 中存在最小数 $M$，即 $M=\sup E$。

\subsection{几个常用不等式}

下面介绍几个常用的不等式。实数 $x$ 的绝对值定义为
\[
  |x|=\begin{cases}
    x, & x\ge0,\\
    -x, & x<0.
  \end{cases}
\]
由此可知，对任何 $x,y\in\mathbb{R}$，有
\[
  -|x|\le x\le |x|,\qquad -|y|\le y\le |y|,
\]
将此两个不等式相加，即得
\[
  -(|x|+|y|)\le x+y\le |x|+|y|,
\]
因此有
\[
  |x+y|\le |x|+|y|.
\]
这一不等式通常称做\impo{三角不等式}。运用三角不等式，可以得到
\[
  \bigl||x|-|y|\bigr|\le |x-y|.
\]

\noindent\impo{伯努利 (Bernoulli) 不等式}\quad 对任意的 $x\ge -1$ 和任意正整数 $n$，有
\[
  (1+x)^n\ge 1+nx.
\]

\noindent\impo{证明}\quad 我们采用数学归纳法来证明这一不等式。当 $n=1$ 时，上述不等式显然成立。假设我们已经证明了
\[
  (1+x)^{n-1}\ge 1+(n-1)x
\]
对一切的 $x\ge -1$ 成立，则有
\[
\begin{aligned}
  (1+x)^n
  &= (1+x)^{n-1}(1+x)\\
  &\ge (1+(n-1)x)(1+x)\\
  &=1+nx+(n-1)x^2\\
  &\ge 1+nx.
\end{aligned}
\]
由数学归纳法原理知，伯努利不等式对一切的正整数 $n$ 成立。

\noindent\impo{算术--几何平均不等式}\quad 对任意 $n$ 个非负实数 $x_1,x_2,\cdots,x_n$，有
\[
  \frac{x_1+x_2+\cdots+x_n}{n}\ge \sqrt[n]{x_1x_2\cdots x_n}.
\]

\noindent\impo{证明}\quad 我们同样采用数学归纳法来证明这一不等式。当 $n=1$ 时，上述不等式显然成立。假设我们已经证明了，对任意 $n-1$ 个非负实数，算术--几何平均不等式成立。下面我们来考虑任意 $n$ 个非负实数 $x_1,x_2,\cdots,x_n$ 的情形。不妨假定 $x_n$ 是这 $n$ 个数中最大的一个，并且令
\[
  y=\frac{x_1+x_2+\cdots+x_{n-1}}{n-1},
\]
则有
\[
  x_n\ge y=\frac{x_1+x_2+\cdots+x_{n-1}}{n-1}
  \ge \sqrt[n-1]{x_1x_2\cdots x_{n-1}},
\]
从而有
\[
\begin{aligned}
  \left(\frac{x_1+x_2+\cdots+x_n}{n}\right)^n
  &=\left(y+\frac{x_n-y}{n}\right)^n\\
  &=y^n+ny^{n-1}\frac{x_n-y}{n}+\cdots+\left(\frac{x_n-y}{n}\right)^n\\
  &\ge y^n+ny^{n-1}\frac{x_n-y}{n}\\
  &=y^{n-1}x_n\\
  &\ge x_1x_2\cdots x_n,
\end{aligned}
\]
即有
\[
  \frac{x_1+x_2+\cdots+x_n}{n}\ge \sqrt[n]{x_1x_2\cdots x_n}.
\]
由数学归纳法原理知，算术--几何平均不等式对一切的正整数成立。

\subsection{常用记号}

全体正整数组成的集合、全体整数构成的集合、全体有理数组成的集合和全体实数组成的集合是本书最常遇到的数集，我们约定分别用空心字母 $\mathbb{N},\mathbb{Z},\mathbb{Q}$ 和 $\mathbb{R}$ 表示，即
\[
\begin{array}{ll}
  \mathbb{N} & \text{表示全体正整数组成的集合；}\\
  \mathbb{Z} & \text{表示全体整数组成的集合；}\\
  \mathbb{Q} & \text{表示全体有理数组成的集合；}\\
  \mathbb{R} & \text{表示全体实数组成的集合.}
\end{array}
\]
显然有
\[
  \mathbb{N}\subset\mathbb{Z}\subset\mathbb{Q}\subset\mathbb{R}.
\]
本书中常要用到以下形式的实数子集：
\[
  [a,b]=\{x\in\mathbb{R}:a\le x\le b\};
\]
\[
  (a,b)=\{x\in\mathbb{R}:a<x<b\};
\]
\[
  (a,b]=\{x\in\mathbb{R}:a<x\le b\};
\]
\[
  [a,b)=\{x\in\mathbb{R}:a\le x<b\}.
\]
这里 $a,b\in\mathbb{R}$，且 $a<b$。

推广区间的概念，把
\[
  (-\infty,b),\quad (-\infty,b],\quad (a,+\infty),\quad [a,+\infty),\quad (-\infty,+\infty)
\]
都称做\impo{无穷区间}。读者应该清楚知道上面每个无穷区间的精确定义，例如 $(-\infty,b]=\{x\in\mathbb{R}:x\le b\}$。

今后，我们提及的区间必是以上形式的区间之一。特别地，对 $a\in\mathbb{R}$，我们把开区间 $(a-\varepsilon,a+\varepsilon)(\varepsilon>0)$ 称为 $a$ 的一个 $\varepsilon$\impo{邻域}，记为 $U(a,\varepsilon)$；而把 $(a-\varepsilon,a+\varepsilon)\setminus\{a\}$ 称为 $a$ 的一个\impo{去心邻域}，记为 $U_0(a,\varepsilon)$。

在本书第一册中，我们还要遇到平面内的点集。在平面直角坐标系中，平面上的点集 $E$ 可以通过描述 $E$ 中元素的坐标所具有的特性来定义。如 $E=\{(x,y):x^2+y^2=1\}$ 就是平面内以原点 $O$ 为圆心，以 $1$ 为半径的圆周上的点组成的集合。

\section{函数的概念}

\subsection{函数的定义}

当我们观察或者研究客观世界的各种事物时，会遇到很多不同的量。这些量一般来说都是在不断变化的，这是客观世界不断变化、不断运动、不断发展在量的方面的体现。例如，当我们研究一个自由落体运动时，就会遇到地球的引力、空气的浮力、时间、物体的高度等各种量，这些量都是变化的，叫做变量。我们知道，客观世界中各个事物之间是相互联系、相互依赖与相互制约的。因此，从量的方面来看，各个变量之间也是相互联系和相互制约的。变量之间的这种相互依赖的关系就是所谓的函数关系。

例如，若在研究自由落体运动时设初始高度为 $h_0$，则在不计空气阻力等其他因素的前提下，该物体距地面的高度 $h$ 与时间 $t$ 有以下关系：
\begin{equation}
  h=h_0-\frac12gt^2,\qquad t\in[0,\sqrt{h_0/g}],
  \tag{1.2.1}
\end{equation}
其中 $g$ 是重力加速度。

在以上自由落体运动过程中，$g$ 始终保持不变，因此它是一个常量；当时间 $t$ 在 $[0,\sqrt{h_0/g}]$ 中不断变化时，$h$ 也随着 $t$ 的变化而不断地变化，因此它们都是变量。在这一过程中，只要知道时间 $t$，根据关系 (1.2.1)，我们就可以知道物体距地面的高度 $h$。换句话说，对于每一个 $t\in[0,\sqrt{h_0/g}]$，根据 (1.2.1) 所规定的对应法则，就有唯一确定的一个高度 $h$ 与之对应。这就是自由落体运动中的函数关系。一般地，我们有

\noindent\impo{定义 1.2.1}\quad 对于给定的集合 $X\subseteq\mathbb{R}$，如果存在某种对应法则 $f$，使得对 $X$ 中的每一个数 $x$，在 $\mathbb{R}$ 中存在唯一的数 $y$ 与之对应，则称对应法则 $f$ 为从 $X$ 到 $\mathbb{R}$ 的一个\impo{函数}，记做
\[
  f:X\to\mathbb{R},
\]
\[
  x\mapsto y=f(x),
\]
其中 $y$ 称为 $f$ 在点 $x$ 的值，$X$ 称为函数 $f$ 的定义域，数集 $\{f(x):x\in X\}$ 称为函数 $f$ 的值域，记为 $f(X)$；$x$ 称做自变量，$y$ 称做因变量。

简言之，函数就是由定义域到它的值域的一个单值对应。此外，构成一个函数必须具备三个基本要素：定义域、值域和对应法则。但其中最重要的两个基本的要素是定义域和对应法则，而值域可由定义域和对应法则所确定。通常，我们用 $y=f(x)(x\in X)$ 记定义在 $X$ 上的一个函数。有时，我们也用 $y=g(x),y=\varphi(x)$ 等记号来表示不同的函数。

在自由落体的运动中，$X=[0,\sqrt{h_0/g}]$ 是定义域，对应法则是 $f$：对每一个 $t\in[0,\sqrt{h_0/g}]$，由解析式 $h=h_0-\frac12gt^2$ 来确定对应的数 $h\in\mathbb{R}$。由于对每一个 $t$，此解析式可确定唯一的数 (即物体高度) 与之对应。因此，物体自由下落这个过程确定了其高度与时间的函数关系。

设 $y=f(x)$ 是定义在 $X$ 上的一个函数，通常称平面点集
\[
  T=\{(x,y):y=f(x),x\in X\}
\]
为该函数的\impo{图像}。一般来讲，它是平面内的一条曲线，它和任何一条平行于 $y$ 轴的直线至多只有一个交点。

在中学数学中，我们所遇到的主要是如下的六类函数：
\begin{enumerate}
  \item 常值函数 $y=C$；
  \item 幂函数 $y=x^\alpha\ (\alpha>0)$；
  \item 指数函数 $y=a^x\ (a>0)$；
  \item 对数函数 $y=\log_a x\ (a>0,\ a\ne1)$；
  \item 三角函数 $y=\sin x,\ y=\cos x,\ y=\tan x,\ y=\cot x$；
  \item 反三角函数 $y=\arcsin x,\ y=\arccos x,\ y=\arctan x,\ y=\operatorname{arccot} x$。
\end{enumerate}
这些函数统称为\impo{基本初等函数}，它们的定义域是使得其解析式有意义的数 $x$ 的全体。在本书中，它们仍将是我们研究的主要对象。基本初等函数的主要特点是函数关系有明显的解析表达式。值得注意的是，这些解析式有着各自特殊的意义。如 $y=\sin x$，我们是用以下方式来确定这个函数关系的。

对每个 $x\in[0,2\pi)$，在 $O\xi\eta$ 坐标平面中取射线 $\overrightarrow{OP}$，使得 $O\xi$ 轴正向逆时针旋转到 $\overrightarrow{OP}$ 的角度为 $x$。设 $\overrightarrow{OP}$ 与单位圆周的交点为 $P$，则 $P$ 的纵坐标即是 $y=\sin x$ 的值 (图 1.2.1)；再规定 $x$ 与 $x+2k\pi\ (k\in\mathbb{Z})$ 取相同的函数值，就得到了定义在 $(-\infty,+\infty)$ 上的函数 $y=\sin x$。

另外，在中学数学中，对某些函数关系也没有给出精确定义。如，在初等数学中就很难弄清楚 $(\sqrt2)^\pi$ (它是函数 $f(x)=x^\pi$ 在 $x=\sqrt2$ 处的取值) 是如何定义的。我们将在后续章节中给出它们的精确定义。今后对于具有解析表达式的函数，若无特殊说明，该函数的定义域即为使得该解析式有意义的自变量的全体所构成的集合。因此，有时候我们对具有解析式的函数不特别指出其定义域。

作为练习，请读者做出基本初等函数的图像。

函数的表示法有穷举法、描述法、列表法及图形法等，以上表示法在中学教科书中已有相应的例子，我们在这里不再赘述。下面我们给出几个中学数学中未曾出现的但在数学分析中经常用到的函数的例子。

\noindent\impo{例 1.2.1 (符号函数)}\quad 所谓符号函数，是指如下函数：
\[
  y=\operatorname{sgn}x=
  \begin{cases}
    1, & x>0,\\
    0, & x=0,\\
    -1, & x<0,
  \end{cases}
\]
其图像如图 1.2.2 所示。这个函数的表达式是分段给出的，通常称这样的函数为分段函数。读者应该注意它是一个函数，而不是三个函数。

\noindent\impo{例 1.2.2 (狄利克雷 (Dirichlet) 函数)}\quad 函数
\[
  y=D(x)=
  \begin{cases}
    1, & x\in\mathbb{Q},\\
    0, & x\in\mathbb{R}\setminus\mathbb{Q}
  \end{cases}
\]
称为狄利克雷函数。这一函数的图像，是分布在 $y=0$ 和 $y=1$ 这两条直线上的两个离散的点集 (参见图 1.2.3)。注意，这样函数的图像是无法精确画出的。

\noindent\impo{例 1.2.3 (高斯 (Gauss) 取整函数)}\quad 取整函数定义为 $y=[x]$，其中 $[x]$ 表示不超过 $x$ 的最大整数，即若 $n\le x<n+1$，则 $[x]=n$。其图像如图 1.2.4 所示。

\noindent\impo{例 1.2.4 (黎曼 (Riemann) 函数)}\quad 所谓的黎曼函数，是指由下式确定的函数：
\[
  y=R(x)=
  \begin{cases}
    \dfrac1p, & x=\dfrac qp\in(0,1),\ p,q\text{ 为互素的正整数},\\
    0, & x\in[0,1]\setminus\mathbb{Q},\\
    1, & x=0\text{ 或 }1.
  \end{cases}
\]
它的图像也不是连续不断的曲线 (参见图 1.2.5)，也是无法精确画出的。

\noindent\impo{例 1.2.5 (特征函数)}\quad 设 $E\subset\mathbb{R}$，定义在 $\mathbb{R}$ 中的函数
\[
  y=\chi_E(x)=
  \begin{cases}
    1, & x\in E,\\
    0, & x\notin E
  \end{cases}
\]
称为集 $E$ 的特征函数。

\subsection{由已知函数构造新函数的方法}

从已知函数出发，可以通过代数四则运算、限制与延拓、复合运算和取反函数等手段构造出许多新的函数。

\noindent\impo{1。函数的四则运算}

设 $y=f_j(x),x\in X_j\subset\mathbb{R}\ (j=1,2)$ 为两个已知函数，且 $X=X_1\cap X_2\ne\varnothing$，则我们可以利用实数的四则运算构造新函数如下：
\[
  (f_1\pm f_2)(x)=f_1(x)\pm f_2(x),\qquad x\in X;
\]
\[
  (f_1f_2)(x)=f_1(x)f_2(x),\qquad x\in X;
\]
\[
  \frac{f_1}{f_2}(x)=\frac{f_1(x)}{f_2(x)}\ (f_2(x)\ne0),\qquad x\in X.
\]
显然，函数的加法、减法与乘法运算可推广到任意有限个函数的情形，而且对于加法和乘法运算，它们具有交换律与结合律。

\noindent\impo{2。函数的限制与延拓}

设函数 $f(x),x\in X_1$ 和 $g(x),x\in X_2$ 满足：$X_1\subset X_2$，且 $f(x)=g(x),\forall x\in X_1$，则称 $f(x)$ 是 $g(x)$ 在 $X_1$ 上的\impo{限制}，而 $g(x)$ 是 $f(x)$ 在 $X_2$ 上的\impo{延拓}。这样，我们就可以从一个已知函数出发，通过限制和延拓来产生新函数。

这里须特别指出的是，对于两个函数而言，只有当它们的定义域及对应关系完全一样时，才能说它们是相等的！例如，$y_1=|x|,x\in(-\infty,+\infty)$ 和 $y_2=x\operatorname{sgn}x,x\in(-\infty,+\infty)$，虽然形式不同，但由于它们具有同样的定义域和对应关系，所以 $y_1$ 和 $y_2$ 是同一个函数；而 $f(x)=x^2,x\in(-\infty,+\infty)$ 和 $g(x)=x^2,x\in(0,+\infty)$ 就是两个不同的函数。

\noindent\impo{3。函数的复合}

函数的复合运算是数学中一种重要的运算。设 $y=f_j(x),x\in X_j\subset\mathbb{R}\ (j=1,2)$ 为两个函数，若 $Y_1=f_1(X_1)\subseteq X_2$，则定义在 $X_1$ 上的函数 $y=f_2(f_1(x))$ 称为 $f_1$ 和 $f_2$ 的\impo{复合函数}，记做 $f_2\circ f_1:X_1\to\mathbb{R}$。通常称 $f_1$ 为该复合函数的\impo{内函数}，$f_2$ 为\impo{外函数}。

同样我们可以考虑多个函数的复合 (如果复合有意义的话)。一般来说，函数复合不具有交换律。例如，取 $f(x)=x^2+1,x\in\mathbb{R},g(x)=x^4,x\in\mathbb{R}$，则 $f(g(x))=x^8+1,x\in\mathbb{R}$，而 $g(f(x))=(x^2+1)^4,x\in\mathbb{R}$。

读者要特别注意是：函数的复合运算可以进行的前提条件是，外函数的定义域必须包含内函数的值域。当然，当两个函数的定义域都是 $\mathbb{R}$ 时，则复合运算可以自由地进行。

函数的复合运算是构造新函数的重要途径。

\noindent\impo{例 1.2.6}\quad 设 $f(x)=\dfrac{x}{\sqrt{1+x^2}}$，试求出 $n$ 次复合函数 $\underbrace{f\circ f\circ\cdots\circ f}_{n\text{ 个}}(x)$ 的表达式。

\noindent\impo{解}\quad 显然，$f(x)$ 的定义域是 $(-\infty,+\infty)$，因此所要讨论的复合运算是可行的。从
\[
  f\circ f(x)=
  \frac{\dfrac{x}{\sqrt{1+x^2}}}{\sqrt{1+\dfrac{x^2}{1+x^2}}}
  =\frac{x}{\sqrt{1+2x^2}}
\]
我们猜出 $n$ 次复合的结果应该为
\[
  \underbrace{f\circ f\circ\cdots\circ f}_{n\text{ 个}}(x)=\frac{x}{\sqrt{1+nx^2}}.
\]
下面用归纳法证之。当 $n=1$ 时，结论显然成立。现在假定当 $n=k$ 时，有
\[
  \underbrace{f\circ f\circ\cdots\circ f}_{k\text{ 个}}(x)=\frac{x}{\sqrt{1+kx^2}},
\]
则当 $n=k+1$ 时，有
\[
  f\circ(\underbrace{f\circ\cdots\circ f}_{k\text{ 个}}\circ f)(x)
  =f\left(\frac{x}{\sqrt{1+kx^2}}\right)
  =\frac{x}{\sqrt{1+(k+1)x^2}}.
\]
这样，由数学归纳法原理知，对一切的 $n$，有
\[
  \underbrace{f\circ f\circ\cdots\circ f}_{n\text{ 个}}(x)=\frac{x}{\sqrt{1+nx^2}}.
\]

\noindent\impo{4。反函数}

在函数的定义中，我们要求对 $\forall x\in X$，在 $\mathbb{R}$ 中必须有唯一的数 $f(x)$ 与之对应，但我们并没有要求 $X$ 中的任何两个不同的数 $x_1$ 和 $x_2$，在 $\mathbb{R}$ 中与之对应的两个数 $f(x_1)$ 和 $f(x_2)$ 也必须不同。最极端的例子是常值函数，其所有点上的函数值取同一个数。

设 $f:X\to Y$ 是一个函数。若对任意的 $x_1,x_2\in X$，只要 $x_1\ne x_2$，就有 $f(x_1)\ne f(x_2)$ 成立，则称 $f(x)$ 是\impo{单的}；若 $Y=f(X)$，则称 $f(x)$ 为\impo{满的}；若 $f(x)$ 既是单的又是满的，则称它为一个\impo{一一对应}。

\noindent\impo{定义 1.2.2}\quad 设 $f:X\to Y$ 是一个一一对应。定义函数 $g:Y\to X$ 如下：对任意的 $y\in Y$，函数值 $g(y)$ 规定为由关系式 $y=f(x)$ 所唯一确定的 $x\in X$。这样定义的函数 $g(y)$ 称为是函数 $f(x)$ 的\impo{反函数}，记为 $g=f^{-1}$。

反函数 $x=f^{-1}(y)$ 的定义域和值域恰为原来函数 $y=f(x)$ 的值域和定义域。函数 $f$ 和其反函数 $f^{-1}$ 满足：
\[
  f(f^{-1}(y))=y,\qquad \forall y\in Y,
\]
\[
  f^{-1}(f(x))=x,\qquad \forall x\in X.
\]
习惯上，我们总是把自变量记为 $x$，因变量记为 $y$。因此，为了遵从统一性，$y=f(x)$ 的反函数也记为 $y=f^{-1}(x)$。这样一来，从几何上来看，$y=f(x)$ 的图像与它的反函数 $y=f^{-1}(x)$ 的图像正好关于直线 $y=x$ 对称。

在中学数学中，我们曾经学习过反三角函数。回想一下，在定义这类函数时，我们必须将三角函数限制在某个区间上。这说明，一个函数在它的定义域中可能不是单的，但把它做适当的限制后则常常可以变成单的，从而可以定义它的反函数。

\noindent\impo{例 1.2.7}\quad 求 $y=f(x)=\dfrac{e^x-e^{-x}}{2}$ 的反函数，这里 $e=2.7182\cdots$ 为一常数。

\noindent\impo{解}\quad 固定 $y$，为了从方程 $y=\dfrac{e^x-e^{-x}}{2}$ 中求出 $x$，我们令 $e^x=z$，则有
\[
  2zy=z^2-1.
\]
解之，得
\[
  z=y\pm\sqrt{y^2+1}.
\]
由于 $z=e^x>0$，故舍去 $y-\sqrt{y^2+1}$，从而有 $x=\ln(y+\sqrt{y^2+1})$。因此，所求的反函数为
\[
  y=f^{-1}(x)=\ln(x+\sqrt{x^2+1}),\qquad x\in(-\infty,+\infty).
\]

本题中的函数 $f(x)$ 称为双曲正弦函数，并记为 $f(x)=\sinh x$ (以后我们将给出所有双曲函数的定义)，它的反函数 $\sinh^{-1}x$ 称为反双曲正弦函数。双曲正弦函数和反双曲正弦函数的图像如图 1.2.6 所示。

\noindent\impo{例 1.2.8}\quad 设 $y=f(x)$ 为定义在 $X$ 上的一个函数，并且记 $Y=f(X)$。证明：若存在 $Y$ 上定义的函数 $g(y)$，使得 $g(f(x))=x$，则 $f(x)$ 的反函数存在，而且 $g=f^{-1}$。

\noindent\impo{证明}\quad 容易验证，复合函数是单的充分必要条件是其内、外函数都是单的。由 $g(f(x))=x$ 是 $X$ 的恒等函数可知，$f(x)$ 是单的，从而 $f(x)$ 是 $X$ 到 $Y$ 的一一对应。于是，$f(x)$ 的反函数存在，而且对 $\forall y\in Y=f(X)$，有 $g(y)=g(f(x))=x=f^{-1}(y)$，即有 $g=f^{-1}$。

作为本节的结束，我们来讨论一下分式线性函数的反函数和复合函数。形如
\[
  y=\frac{ax+b}{cx+d}\qquad (ad-bc\ne0)
\]
的函数称为\impo{分式线性函数}，这里的条件 $ad-bc\ne0$ 是防止它变成一个常值函数。

由 $y=\dfrac{ax+b}{cx+d}$，可推出 $x=\dfrac{b-dy}{cy-a}$，而且有
\[
  (-d)(-a)-bc=ad-bc\ne0.
\]
这表明，任何一个分式线性函数都具有反函数，而且其反函数仍然为分式线性函数。

另外，设
\[
  f_j(x)=\frac{a_jx+b_j}{c_jx+d_j}\quad (a_jd_j-b_jc_j\ne0),\qquad j=1,2,
\]
则容易导出
\[
  f_1\circ f_2(x)=
  \frac{(a_1a_2+b_1c_2)x+a_1b_2+b_1d_2}
  {(c_1a_2+d_1c_2)x+c_1b_2+d_1d_2},
\]
而且
\[
\begin{aligned}
 &(a_1a_2+b_1c_2)(c_1b_2+d_1d_2)
 -(a_1b_2+b_1d_2)(c_1a_2+d_1c_2)\\
 &\qquad=(a_1d_1-b_1c_1)(a_2d_2-b_2c_2)\ne0.
\end{aligned}
\]
这表明，任意两个分式线性函数的复合还是分式线性函数。

分式线性函数在数学的许多分支中起着重要的作用。这个函数类有着许多有趣的性质。例如，每个分式线性函数 $y=\dfrac{ax+b}{cx+d}(ad-bc\ne0)$ 可对应一个 $2\times2$ 矩阵 $\begin{pmatrix}a&b\\ c&d\end{pmatrix}$。有趣的是，任何分式线性函数的反函数所对应的矩阵刚好是它对应矩阵的逆乘以它的行列式的 $-1$ 倍，而任意两个分式线性函数的复合函数所对应的矩阵则刚好是它们各自对应矩阵的乘积.

细心的读者可以发现，对 $y=f(x)=\dfrac{ax+b}{cx+d}$ 来说，当 $c\ne0$ 时，$f(x)$ 在 $x=-\dfrac dc$ 没有定义；另外 $f(x)$ 的值域也不包含 $\dfrac ac$。如果我们形式地令 $f(-\dfrac dc)=\infty$ 和 $f(\infty)=\dfrac ac$，而当 $c=0$ 时令 $f(\infty)=\infty$，则可以验证 $f(x)$ 是 $\mathbb{R}\cup\{\infty\}$ 到 $\mathbb{R}\cup\{\infty\}$ 的一个一一对应，从而以上的求反函数和复合的运算可以自由地进行。

\section{函数的性质}

在本节中，我们主要介绍某些函数具有的一些特别性质。这些性质的讨论是我们以后经常要遇到的，因此要求读者能用准确的数学语言来叙述它们的定义和一些等价的命题。

\subsection{函数的有界性}

设 $y=f(x)$ 是定义在 $X$ 上的函数。若存在常数 $M$，使得对 $\forall x\in X$，都有 $f(x)\le M$，则称 $f(x)$ 在 $X$ 上有\impo{上界}，同时称 $M$ 是 $f(x)$ 的一个上界；若存在常数 $m$，使得对 $\forall x\in X$，都有 $f(x)\ge m$，则称 $f(x)$ 在 $X$ 上有\impo{下界}，同时称 $m$ 是 $f(x)$ 的一个下界；若 $f(x)$ 在 $X$ 上既有上界 $M$ 又有下界 $m$，则称 $f(x)$ 在 $X$ 上\impo{有界}。此时，对 $\forall x\in X$，有 $m\le f(x)\le M$。若 $|f(x)|\le M(x\in X)$，则称 $M$ 是 $f(x)$ 的一个\impo{界}。

容易看出，$f(x)$ 在 $X$ 上有界的充分必要条件是存在 $M>0$，使得当 $x\in X$ 时，有 $|f(x)|\le M$，即 $f(X)\subseteq[-M,M]$。换句话说，$f(x)$ 在 $X$ 有界也就是其值域 $f(X)$ 是一个有界集。

从几何上来看，若函数 $y=f(x)$ 在 $X$ 上有上界 $M$，则 $f(x)$ 的图像将位于直线 $y=M$ 的下面；若 $f(x)$ 有下界 $m$，则 $f(x)$ 的图像在直线 $y=m$ 的上面；若 $f(x)$ 有界，则必存在 $M>0$，使得 $f(x)$ 的图像位于直线 $y=M$ 和 $y=-M$ 之间。

\noindent\impo{例 1.3.1}\quad 函数 $y=\sin x$ 和 $y=\cos x$ 在 $\mathbb{R}$ 上有界，$M=1$ 就是它们的一个界。而对函数 $y=x^n,\ x\in(-\infty,+\infty)$ ($n$ 为正整数) 来说，当 $n$ 为奇数时，它既无上界，又无下界；而当 $n$ 为偶数时，它无上界，但有下界 $m=0$。

若 $f(x)$ 不是 $X$ 上的有界函数，则称该函数在 $X$ 上\impo{无界}。如何用肯定的语气来叙述函数的无界性呢？我们分析一下，$f(x)$ 在 $X$ 上无界的等价说法是，任何正数 $M$ 都不是 $y=f(x)$ 在 $X$ 上的界。而 $M$ 不是 $y=f(x)$ 在 $X$ 上的界，就等价于，必然 $\exists x_0\in X$，使得 $|f(x_0)|>M$。因此，若用肯定的语气，函数 $f(x)$ 在 $X$ 上无界可以叙述为：$\forall M>0$，$\exists x_0\in X$，使得 $|f(x_0)|>M$。

同理我们也可以给出函数 $f(x)$ 在 $X$ 上无上界或无下界的定义，请读者自己给出这些定义的精确数学语言的描述。

\noindent\impo{例 1.3.2}\quad 证明：$f(x)=\dfrac{\sin\frac1x}{x}$ 在 $(0,1]$ 区间上既无上界也无下界，但它却在 $[a,1]$ 上有界，这里 $1>a>0$。

\noindent\impo{证明}\quad $\forall M>0$，取正整数 $n_0>M$ 并取
\[
  x_0=\frac{1}{2n_0\pi+\pi/2}\in(0,1],
\]
则有
\[
  f(x_0)=\frac{1}{x_0}=2n_0\pi+\frac{\pi}{2}>n_0>M.
\]
这就证明了 $y=\dfrac{\sin\frac1x}{x}$ 在 $(0,1]$ 上是无上界的。

取
\[
  x_1=\frac{1}{2n_0\pi-\pi/2}\in(0,1],
\]
则有
\[
  f(x_1)=\frac{-1}{x_1}=-\left(2n_0\pi-\frac{\pi}{2}\right)<-n_0<-M.
\]
这就证明了 $y=\dfrac{\sin\frac1x}{x}$ 在 $(0,1]$ 上是无下界的。

对 $a>0$，取 $M=\dfrac1a>0$，则 $\forall x\in[a,1]$，有
\[
  \left|\frac{\sin\frac1x}{x}\right|\le \left|\frac1x\right|\le M,
\]
即 $y=\dfrac{\sin\frac1x}{x}$ 在 $[a,1]$ 上有界。

\subsection{函数的单调性}

设 $y=f(x)$ 是定义在 $X$ 上的一个函数。若对任意的 $x_1,x_2\in X$，只要 $x_1<x_2$，便有 $f(x_1)\le f(x_2)(f(x_1)\ge f(x_2))$，则称 $f(x)$ 在 $X$ 上是\impo{单调上升 (下降) 函数}或\impo{单调递增 (递减) 函数}。在上述不等式中将``$\le$'' (``$\ge$'') 换成``$<$'' (``$>$'')，则称 $f(x)$ 在 $X$ 上是\impo{严格单调上升 (下降) 函数}。单调上升函数和单调下降函数统称为\impo{单调函数}。

严格单调的函数是定义域到值域的一一对应，所以它存在反函数。但存在反函数的函数未必是单调的。

\noindent\impo{例 1.3.3}\quad $y=\sin x(x\in\mathbb{R})$ 不是单调函数，但当其定义域 $X$ 取为 $\left[-\dfrac\pi2,\dfrac\pi2\right]$ 时，它是严格单调递增函数；当 $X$ 取为 $\left[\dfrac\pi2,\dfrac{3\pi}2\right]$ 时，它是严格单调递减函数。因此该函数在这两个区间都分别存在反函数。在习惯上，我们用 $y=\arcsin x$ 表示 $y=\sin x$ 在 $\left[-\dfrac\pi2,\dfrac\pi2\right]$ 上的反函数。

\noindent\impo{例 1.3.4}\quad 考查函数
\[
  y=f(x)=
  \begin{cases}
    x, & x\in[0,1]\cap\mathbb{Q},\\
    1-x, & x\in[0,1]\setminus\mathbb{Q}.
  \end{cases}
\]
容易看出，该函数的反函数就是它自己，但 $f(x)$ 在 $[0,1]$ 的任何子区间都不是单调的。

\subsection{函数的周期性}

设 $y=f(x)$ 是在 $X$ 上有定义的函数。若存在 $T>0$，使得对任意 $x\in X$ 有 $f(x+T)=f(x)$，则称 $f(x)$ 为\impo{周期函数}，$T$ 称为 $f(x)$ 的一个\impo{周期}。若存在一个最小的周期 $T_0$，则称 $T_0$ 为 $f(x)$ 的\impo{基本周期}。

显然，以 $T$ 为周期的周期函数 $f(x)$ 的定义域 $X$ 必须满足条件：对一切的 $x\in X$，有 $T+x\in X$。

\noindent\impo{例 1.3.5}\quad 函数 $y=\tan x$ 的定义域为
\[
  \mathbb{R}\setminus\left\{\frac\pi2+n\pi:n\in\mathbb{Z}\right\},
\]
它是以 $\pi$ 为基本周期的周期函数。

这里须特别指出的是，并非所有的周期函数都有基本周期。

\noindent\impo{例 1.3.6}\quad 容易验证，任何正有理数都是狄利克雷函数 $D(x)$ (参看例 1.2.2) 的周期。因此，它没有基本周期。

若一个定义在闭区间 $[a,b]$ 上的函数 $f(x)$ 满足 $f(a)=f(b)$，则可将它延拓成 $(-\infty,+\infty)$ 上以 $T=b-a$ 为周期的周期函数 $\tilde f(x)$。

\noindent\impo{例 1.3.7}\quad 设 $f(x)$ 为定义在 $\mathbb{R}$ 上的周期函数，并且有基本周期 $\tau>0$。证明：若 $\forall x\in(0,\tau)$，有 $f(x)\ne f(0)$，则 $g(x)=f(x^2)$ 不是周期函数。

\noindent\impo{证明}\quad 用反证法。假定 $g(x)$ 是周期函数，而且 $T>0$ 是它的一个周期，则有
\begin{equation}
  f((x+T)^2)=g(x+T)=g(x)=f(x^2),\qquad \forall x\in\mathbb{R}.
  \tag{1.3.1}
\end{equation}
以 $x=0$ 代入 (1.3.1) 式，得 $f(T^2)=f(0)$。而 $f(x)$ 是以 $\tau$ 为基本周期的周期函数，因此必有某一正整数 $k$，使得 $T^2=k\tau$，即 $T=\sqrt{k\tau}$。

再以 $x=\sqrt{(k+1)\tau}$ 代入 (1.3.1) 式，得
\[
  f\left(\left(\sqrt{(k+1)\tau}+\sqrt{k\tau}\right)^2\right)=f((k+1)\tau)=f(0),
\]
而
\[
  f\left(\left(\sqrt{(k+1)\tau}+\sqrt{k\tau}\right)^2\right)
  =f\left((2k+1)\tau+2\sqrt{(k+1)k}\tau\right)
  =f(2\sqrt{(k+1)k}\tau),
\]
所以
\[
  f(2\sqrt{(k+1)k}\tau)=f(0).
\]
这又蕴涵着，必有某一正整数 $n$，使得 $2\sqrt{(k+1)k}\tau=n\tau$。由此可得 $\left(\dfrac n2\right)^2=k(k+1)$。这一等式表明，$\dfrac n2$ 是正整数，而且满足 $k<\dfrac n2<k+1$。这显然是不可能的。这一矛盾表明，$g(x)$ 不可能是周期函数。

\subsection{函数的奇偶性}

设 $y=f(x)$ 是定义在 $X$ 上的一个函数，而且 $X$ 是关于原点对称的，即 $x\in X$ 蕴涵着 $-x\in X$。若 $f(x)=-f(-x)$ 对一切的 $x\in X$ 成立，则称 $f(x)$ 是 $X$ 上的\impo{奇函数}；若 $f(x)=f(-x)$ 对一切的 $x\in X$ 成立，则称 $f(x)$ 是 $X$ 上的\impo{偶函数}。

容易看出，偶函数的图像是关于 $y$ 轴对称的，而奇函数的图像是关于原点对称的。

\noindent\impo{例 1.3.8}\quad 设 $y=f(x)$ 是 $X$ 上的奇函数且存在反函数。证明：它的反函数也是奇函数。

\noindent\impo{证明}\quad 由奇函数的定义知，对 $\forall x\in X$，有 $f(-x)=-f(x)$。于是
\[
  -f^{-1}(f(x))=-x=f^{-1}(f(-x))=f^{-1}(-f(x)),\qquad \forall x\in X,
\]
即 $\forall y\in f(X)$，有 $-f^{-1}(y)=f^{-1}(-y)$。这就证明了 $f^{-1}$ 也是奇函数。

\noindent\impo{例 1.3.9}\quad 证明 $f(x)=\ln(x+\sqrt{1+x^2})$ 是奇函数。

\noindent\impo{证明}\quad 由于 $f(x)$ 是奇函数 $y=\sinh x$ 的反函数，故由上例知，该函数是奇函数。

\section{初等函数}

在 §1.2 中已经指出，常值函数、幂函数、指数函数、对数函数、三角函数和反三角函数统称为基本初等函数。从这些基本初等函数出发，经过有限多次加、减、乘、除和复合运算所能得到的所有函数统称为\impo{初等函数}。对基本初等函数及其基本性质，大家已经十分熟悉，这里我们就不再赘述。下面我们只简要地介绍一下双曲函数，这是一类十分有用的初等函数，在以后的学习中常常用到它们。

双曲函数定义如下：
\[
  \sinh x=\frac{e^x-e^{-x}}{2},\qquad
  \cosh x=\frac{e^x+e^{-x}}{2},
\]
\[
  \tanh x=\frac{e^x-e^{-x}}{e^x+e^{-x}},\qquad
  \coth x=\frac{e^x+e^{-x}}{e^x-e^{-x}}.
\]
我们把这四个函数分别称之为\impo{双曲正弦}、\impo{双曲余弦}、\impo{双曲正切}和\impo{双曲余切}。由定义不难看出，$\sinh x,\tanh x,\coth x$ 都是奇函数，而 $\cosh x$ 是偶函数 (参见图 1.4.1 (a),(b))。

由双曲函数定义出发，可以导出许多类似于三角函数的恒等式。例如，有
\[
  \sinh 2x=2\sinh x\cosh x,\qquad
  \cosh^2x-\sinh^2x=1,
\]
\[
  \cosh^2x+\sinh^2x=\cosh 2x,\qquad
  \cosh x=1+2\sinh^2\frac x2.
\]
这些恒等式正好对应于如下的三角函数恒等式：
\[
  \sin 2x=2\sin x\cos x,\qquad
  \cos^2x+\sin^2x=1,
\]
\[
  \cos^2x-\sin^2x=\cos 2x,\qquad
  \cos x=1-2\sin^2\frac x2.
\]
令 $X=\cosh x,\ Y=\sinh x$，则它们满足双曲方程
\[
  X^2-Y^2=1.
\]
这是双曲函数名称由来的原因之一。在单位圆盘上的非欧几何——双曲几何 (俄国人称之为罗巴切夫斯基几何，西方人称之为庞加莱 (Poincaré) 几何) 理论中，双曲函数是基本的研究工具。

\section*{习题一}

\begin{enumerate}
  \item 设集合 $A=\{1,2,4,8\},\ B=\{1,3,6,9\}$，试求 $A\cup B,\ A\cap B$ 和 $A-B$。
  \item 设集合 $A=\{2n-1:n\in\mathbb{N}\},\ B=\{2n:n\in\mathbb{N}\}$，试求 $A\cup B,\ A\cap B,\ A-B$ 和 $B-A$。
  \item $\forall n\in\mathbb{N}$，设 $A_n=\left[-1+\dfrac{1}{2n},1-\dfrac1n\right]$，试求 $\displaystyle\bigcup_{n=1}^{+\infty}A_n$。
\end{enumerate}

\end{document}
